% Standalone report fragment. Requires amsmath, amssymb; no custom macros.
\section{Logical calibration of PA194Q}
\label{sec:q-calibration}

The new machine changes one axiom head and retains the cached-head
commands of the pinned PA194C construction. Its source instruction map
is therefore not asserted to equal PA194C's instruction map. The
preservation argument passes through finite derivations in two
interderivable numerical calculi, then through their complete command
compilers and literal tape simulations.

\subsection{The exact change of axiom}

Put
\[
 p(a,b)=\frac{(a+b)(a+b+1)}2+b,\qquad
 I(a,b)=4p(a,b),\qquad U(i,B)=4p(B,i)+2.
\]
Write $a\to b$ for the numerical constructor $I(a,b)$. In the
separate-diagonal-head calculus let
\[
 P=U(D,A),\qquad Q=U(D,I(A,H)),\qquad R=U(D,H).
\]
The old sixth head and its new replacement are respectively
\[
 L_6(D,A,H)=P\to(Q\to R),\qquad
 L_6^{\mathrm Q}(D,A,H)=Q\to(P\to R).
\tag{Q1}
\]
Every other head, the guarded separation gate, the separate diagonal
head, the numerical constructors, and the inference rules are retained.
There is no change to the numerical parameters $D,A,H$ within any
formula. In particular, (Q1) is an exchange of two premises within one
axiom, not a permutation of the three parameter names or an exchange
of stack positions.

Both postfix emitters have 19 tokens. Writing $A,D,H$ for the macros
that push the old top and load that register, $C$ for the pairing macro,
and $i$ for increment, their words are
\[
\begin{aligned}
 L_6 &: A\,D\,C\,i\,i\;
         A\,H\,C\,D\,C\,i\,i\;
         H\,D\,C\,i\,i\;C\,C,\\
 L_6^{\mathrm Q} &: A\,H\,C\,D\,C\,i\,i\;
         A\,D\,C\,i\,i\;
         H\,D\,C\,i\,i\;C\,C.
\end{aligned}
\tag{Q2}
\]
The exact emitter audit identifies these words inside the complete
190-token AX word, checks both displayed formula blocks in a free algebra,
and verifies that every surrounding token is unchanged. Each block pushes
the incoming top and leaves the older stack intact, so the replacement
preserves every other emitted stack entry.

\subsection{A proof translation independent of quantifier distribution}

The unchanged restricted implicational basis derives, uniformly in
all numerical codes $a,b,c$, the exchange formula
\[
 \bigl(a\to(b\to c)\bigr)\to
                  \bigl(b\to(a\to c)\bigr).
\tag{Q3}
\]
This is already the theorem \texttt{Dominated.exchange} in
\texttt{PA196/Logic/ImplicationBasis.lean}. Its construction uses only
the restricted propositional head, modus ponens, reflexivity of the
numerical bound, and the two inequalities
$a,b\le I(a,b)$. In particular it does not use either version of the
sixth head. It also needs no rule treating code one as falsity.

For clarity, (Q3) is an intuitionistic implicational tautology. From
assumptions $f:a\to(b\to c)$, $b$, and $a$, two applications of modus
ponens give $c$; three applications of the implicational deduction
construction give (Q3). The $K$ and $S$ combinators used by that
construction are derived from the same unchanged restricted basis.
The optional certificate supplied with this report checks a finite
$K,S,\mathrm{MP}$ derivation of this exact exchange formula.

Apply (Q3) at $(a,b,c)=(Q,P,R)$ and use the new axiom
$Q\to(P\to R)$ as its minor premise. This derives the old axiom
$P\to(Q\to R)$. Conversely, (Q3) at $(P,Q,R)$ derives the new axiom
from the old one. These are uniform finite derivations for every
triple of parameter values $D,A,H$.

Let $\mathcal C$ be the pinned separate-diagonal-head calculus and
$\mathcal C_{\mathrm Q}$ the calculus with (Q1)'s replacement.
Transform a $\mathcal C$ derivation recursively: retain every
unchanged axiom; replace a sixth-head leaf by the new sixth-head
axiom followed by the first exchange derivation above; recursively
transform the premises at modus ponens and generalization nodes.
Every conclusion is preserved literally. The converse transformation
uses the other exchange derivation. Thus there are primitive-recursive
proof transformations $F,G$ whose correctness PA proves:
\[
\begin{aligned}
 \operatorname{Prf}_{\mathcal C}(e,B)
   &\Longrightarrow
       \operatorname{Prf}_{\mathcal C_{\mathrm Q}}(F(e),B),\\
 \operatorname{Prf}_{\mathcal C_{\mathrm Q}}(e,B)
   &\Longrightarrow
       \operatorname{Prf}_{\mathcal C}(G(e),B).
\end{aligned}
\tag{Q4}
\]
There is no substitution of arbitrary formulas into a
freshness-restricted gate. Only the numerical parameters of the
unrestricted implicational exchange template are instantiated.
In particular (Q4) proves, inside PA, equivalence of the two
existential refutation predicates with conclusion code one.

\paragraph{The code-zero alias.}
These are numerical derivations, so they include zero among the allowed
formula codes. The identity $I(0,0)=0$ never violates a modus-ponens
step: its code equality remains exactly the equality required by the
rule. In the arithmetic interpretation of formulas, $T(B)$ denotes the
formula produced by a syntactic translation, not an arithmetic truth
predicate. There is a uniform proof-producing construction of the
instancewise PA equivalence
\[
 T(I(a,b))\leftrightarrow(T(a)\to T(b)),
\]
whose exceptional $(a,b)=(0,0)$ case is the fixed equivalence of truth
with truth implying truth. One must not assert literal syntactic
identity in that exceptional case. The new head's three displayed
subformulas are universal codes and hence are nonzero; more generally
the proof translation does not require any nonzero-code assumption.

\subsection{The ten-entry cached command interface}

Let $[a;w]=p(a,w)$ and let $\operatorname{hd},\operatorname{tl}$ be
the total Cantor projections. In particular
$\operatorname{hd}(0)=\operatorname{tl}(0)=0$ and $[0;0]=0$.
The command store is $\langle t,w,d,h\rangle$. The commands retained
from PA194C are
\[
 S\langle t,w,d,h\rangle=
 \begin{cases}
   \mathsf{halt},&h=1,\\
   \langle h,\operatorname{tl}(w),d,\operatorname{hd}(w)\rangle,
       &h\ne1,
 \end{cases}
\]
and, with $y=d+x$,
\[
 \mathsf{MP}_x\langle t,w,d,h\rangle=
 \begin{cases}
   \langle0,w,0,y\rangle,&I(t,y)=h,\\
   \langle0,w,y,0\rangle,&I(t,y)\ne h.
 \end{cases}
\tag{Q5}
\]
Only Shift can halt. The new axiom command, with
$D=d+x$, $H=h+y$, $A=c$, has effect
\[
 \mathsf{AX}^{\mathrm Q}_{c,x,y}\langle t,w,d,h\rangle
  =\langle K(H,D,c),[q_1,\ldots,q_9,t;w],0,0\rangle,
\]
where
\[
 (q_1,\ldots,q_9)=
 (\Gamma(A),L_6^{\mathrm Q},L_5,L_4,L_3,L_2,L_1,
      \rho_8(D,A,H,d),U(d,h)).
\tag{Q6}
\]
Here $K(a,b,c)=G(a,b,c)$ is the unchanged guarded gate defined in the
following all-axiom section, and
$\Gamma(B)=I(B,U(B,B))$. The saved top remains the tenth pushed
entry. All $q_j$ except $q_9$ have residue zero modulo four;
$q_9$ has residue two. The gate has residue zero, including a failed
guard's output zero. Hence none of these generated entries is one.

The general cached-command proof depends on those residues, the ten
positions, and the last two entries. It does not depend on the internal
curried-premise order of $q_2$. For the new AX command define
\[
\begin{aligned}
 \mathsf{Push}&=\mathsf{AX}^{\mathrm Q}_{0,0,0};S^9;
                                      \mathsf{MP}_0,\\
 \mathsf{Pop}&=\mathsf{AX}^{\mathrm Q}_{0,0,0};S^{11},\\
 \mathsf{Discard}&=\mathsf{MP}_1;
     \mathsf{AX}^{\mathrm Q}_{1,0,0};S^{10};\mathsf{MP}_0.
\end{aligned}
\tag{Q7}
\]
These words satisfy, for all displayed natural values,
\[
\begin{aligned}
 \mathsf{Push}\langle t,w,d,h\rangle
       &=\langle0,[t;w],0,0\rangle,\\
 \mathsf{Pop}\langle t,[a;w],d,h\rangle
       &=\begin{cases}\mathsf{halt},&t=1,\\
                         \langle t,w,0,a\rangle,&t\ne1,
          \end{cases}\\
 \mathsf{Discard}\langle t,w,0,0\rangle
       &=\langle0,w,0,0\rangle.
\end{aligned}
\tag{Q8}
\]
Indeed, the first ten Shifts after AX test the cleared cache and the
nine generated entries, and never test the saved top. The next Shift
tests precisely that top. Push's final MP necessarily fails because
an implication has residue zero and the current universal cache has
residue two. Discard's two failed comparisons follow from
$I(t,1)>0$ and $I(U(1,0),0)>0$. These calculations use the total
projection identities, including their zero aliases; they do not
require stacks to be free lists.

For $1\le j\le7$, the head selector remains
\[
 \mathsf{Discard};\mathsf{AX}^{\mathrm Q}_{A,D,H};
                      S^{j+1};\mathsf{Pop}^{9-j}.
\tag{Q9}
\]
It returns $q_j$ on the original arbitrary stack with both caches zero,
starting from any $\langle t,w,0,0\rangle$, even when $t=1$.
In particular the selector for the changed head is $S^3$ followed by
seven Pops after Discard and AX. It returns the new head, not the
old one; the latter is obtained by compiling the derivation in (Q4).
The unchanged generalization, restricted-propositional-head, and
gate selectors have the same proofs and exponents as PA194C.

For completeness, with $\mathsf{Offset}_0$ empty and
$\mathsf{Offset}_i=\mathsf{MP}_i$ for $i>0$, generalization uses
\[
 \mathsf{Gen}_i=\mathsf{Push};\mathsf{Offset}_i;S;
              \mathsf{AX}^{\mathrm Q}_{i,0,0};S^{10}.
\]
It sends $\langle a,w,0,0\rangle$ to
$\langle U(i,a),w,0,0\rangle$ for every $a,w$, including $a=1$.
The MP suffix $S;S;\mathsf{MP}_b;S$ on
\[
          \langle0,[a,I(a,b),0;w],0,0\rangle
\]
halts if $a=1$ or $b=1$, and otherwise returns
$\langle b,w,0,0\rangle$. The dummy zero is essential for preserving
the arbitrary underlying stack. The calculation also covers
$a=b=w=0$ and $I(0,0)=0$.

\subsection{All-store completeness and syntactic soundness}

For a finite $\mathcal C_{\mathrm Q}$ proof $e$, compile axiom leaves
by their selectors, generalization by the preceding word, and an MP
node with major proof $e_1$ of $I(a,b)$ and minor proof $e_0$ of $a$ by
\[
 \mathsf{Discard};\mathsf{Push};W_{e_1};\mathsf{Push};
 W_{e_0};\mathsf{Push};S;S;\mathsf{MP}_b;S.
\]
Induction on the finite proof tree proves the exact invariant:
from every $\langle t,w,0,0\rangle$, the word $W_e$ either halts
or returns $\langle\operatorname{concl}(e),w,0,0\rangle$.
Therefore, if $e$ proves one, the word
\[
       \mathsf{AX}^{\mathrm Q}_{0,0,0};W_e;\mathsf{Pop}
\tag{Q10}
\]
halts from every incoming store. This all-store assertion handles the
inherited scheduler's retention of the store between unsuccessful
trials. An assumption that every trial resets the store would be wrong.
The unchanged exhaustive scheduler eventually tests the word (Q10),
so a refutation yields a source halt.

Conversely annotate the visible top, cached head, and each live
canonical stack entry by derivations in $\mathcal C_{\mathrm Q}$.
Zero has the fixed derivation $\rho_8(0,0,0,0)=0$, so a single
zero proof supplies the total stack's padding. AX emits its own axiom
heads, the saved theorem, and a generalization of the cached theorem;
successful MP applies modus ponens; failed MP returns zero.
Shift moves derived formulas and halts only at a cached proof of one.
This finite proof-store invariant extracts a refutation from every
source halt. The changed head is sound in this invariant because it
is an axiom of the new calculus, and by (Q4) it also has a proof in
the old calculus.

These constructions are primitive recursive. Their correctness uses
induction on finite proofs, total Cantor-stack decoding, natural loop
counters, and finite traces. PA consequently proves
\[
 H_{R_{\mathrm Q}} \leftrightarrow\
   \exists e\,\operatorname{Prf}_{\mathcal C_{\mathrm Q}}(e,1),
 \qquad
 H_{R_{\mathrm C}}\ \leftrightarrow\
   \exists e\,\operatorname{Prf}_{\mathcal C}(e,1).
\tag{Q11}
\]
Here $R_{\mathrm Q}$ is the newly bound source program and
$R_{\mathrm C}$ the pinned cached-head source. Equation (Q4),
not source-step identity, connects their existential halting predicates.

\subsection{Literal machines and the PA biconditional}

The machine-specific certificates bind the literal PA194Q table to
$R_{\mathrm Q}$: every one of its 4,749 source outcomes has a
positive-length, non-premature-halting simulation ending in the exact
intended complete frame, and its verified blank startup reaches the
clean source seed. The corresponding interfaces for PA194C are
retained from the pinned source packet. The generic arithmetic proofs
of the literal transition and register-scan certificates yield
\[
 \mathrm{PA}\vdash H_{M_{\mathrm Q}}\leftrightarrow H_{R_{\mathrm Q}},
 \qquad
 \mathrm{PA}\vdash H_{M_{\mathrm C}}\leftrightarrow H_{R_{\mathrm C}}.
\tag{Q12}
\]
The forward trace construction concatenates the physical simulations
and prepends startup. For the reverse implication, a first literal
halt at time $T$ contradicts $T+1$ nonhalting source steps: their
positive physical durations would give, by determinism, a nonhalting
literal execution extending beyond $T$. This is a bounded arithmetic
argument and does not assume either machine actually runs forever.

Combining (Q4), (Q11), and (Q12) proves the new preservation theorem
\[
 \boxed{\mathrm{PA}\vdash
       H_{M_{\mathrm Q}}\leftrightarrow H_{M_{\mathrm C}}.}
\tag{Q13}
\]
The supplied calibration of $\mathcal C$ gives primitive-recursive
translations between its refutations and PA refutations, with
PA-provable correctness.
Its reverse translation uses arithmetic finite bit comprehension
as a proof-producing induction template, not a uniform truth predicate.
Its forward translation uses the diagonal-freshness replacement of
equality vacuity and the established PA-refutation compiler.
Using that inherited mathematical calibration with (Q13) gives
\[
 \boxed{\mathrm{PA}\vdash
   \bigl(\neg H_{M_{\mathrm Q}}\leftrightarrow
                         \operatorname{Con}(\mathrm{PA})\bigr).}
\tag{Q14}
\]
This is an ordinary mathematical proof with executable finite and
universal symbolic certificates. The new exchange and transport
arguments do not turn the repository's conditional iPRA endpoints
into completed formalizations. No newly compiled end-to-end Lean
or Rocq theorem for (Q14) is asserted.
